% A concrete internal case that turns the proposal's pacing discussion into a
% testable product requirement.  The complete record remains in the appendix.
\chapter{A live lesson in density}
\label{ch:hpo-density}

The manager used the same word for each failure. Three times, while reviewing
a 90-page methods monograph, the manager said the document was
``too dense.'' Three times, the writing agent had already judged the draft
readable. The complaint was right each time, but the repair was different.

That small history is more useful than another universal readability score. It
shows what a reader experiences and what a product has to discover. The first
round concerned the size of the paragraph. The second concerned the number of
new ideas arriving over a page. The third concerned order: the document used
names before it had taught the objects those names referred to.

The figures in this chapter come from one internal engagement, recorded on
24--25 August 2026. They are calibration evidence for a feasibility build,
not estimates of how often technical documents fail and not acceptance
thresholds. The full provenance and the proposed fixture are in
Appendix~\ref{app:hpo-density-record}.

\hvFigureHpoRounds

\section{The first round: one paragraph carried a page}
\label{sec:hpo-packing}

The first version contained a 210-word paragraph. It named six mechanisms and
five citations before the reader had a chance to decide what the paragraph was
trying to establish. Nothing in it was necessarily false. The problem was
that each sentence opened another line of thought while the first one was
still unfinished.

The repair was modest. The author gave each paragraph one job and moved the
qualifications and citations next to the claims they qualified. The mean
length of a prose paragraph fell from 128 words to 90. That count is not a
rule for every document; it is a record of where this reader found relief.

The useful finding would therefore not say ``shorten paragraphs.'' It would
show the source lines, count the new mechanisms and citations in the passage,
and ask whether the paragraph can close one question before it opens another.
The author may split it, add a concrete example, or keep it and explain why
the concentration is necessary.

\begin{cardexample}{The question a packing finding should ask}
\textbf{Observed:} one paragraph introduces six mechanisms and five citations
before it states the consequence.\\
\textbf{Question:} which one result should the reader carry into the next
paragraph?\\
\textbf{Possible repair:} give that result its own paragraph, then unfold the
mechanisms one at a time. The tool does not choose the split.
\end{cardexample}

\section{The second round: the page moved too quickly}
\label{sec:hpo-pacing}

The manager then supplied two documents and said, in effect, ``write at this
pace.'' That instruction was more precise than a target word count. The same
concept detector was run over the survey and the two exemplars. The survey
introduced a new concept about every 70 words. The exemplars introduced one
about every 83 and 253 words. The longer figure belonged to an 800-page
single-project monograph: it was not a universal ideal, but it revealed the
standard this manager had in mind.

The detector also counted paragraphs that introduced at least three new
concepts. The proportions were 12.9 per cent for the survey, 5.9 per cent for
one exemplar, and 0.8 per cent for the other. The absolute counts are noisy;
the comparison is useful because the same imperfect detector was applied to
all three documents.

The repair was not to pad the survey. It was to let one idea stay on stage
long enough to become familiar. A concrete case came first. The prose named
the mechanism. A figure showed it. A short read-back paragraph told the reader
what the figure had changed. Only then did the author introduce the next
distinction.

This changes the product question. A pacing finding may recommend fewer named
things in a unit, not more words in the same unit. It should report the target
document's position relative to its supplied exemplars and list the paragraphs
with the largest bursts. A number without a nearby passage is not an editorial
instruction.

\hvFigureDensityLadder

\section{The third round: names arrived before their objects}
\label{sec:hpo-ordering}

The final complaint was more specific: ``Do not use a concept before you have
explained it.'' A first-use audit found roughly 25 genuine violations. The
running example used the name PPO on its first page and never introduced the
algorithm. Two acquisition-function names appeared a section before the
chapter taught what an acquisition was.

The reader could recognise the words and still be unable to follow the
argument. Recognition is not understanding. A glossary at the end would not
repair the first encounter, because the reader had already had to guess why
the name mattered.

The repair pattern was simple. Start with the ordinary object, then name it:
``a population-based method, BG-PBT, examined in Chapter 6.'' When the order
must be anticipatory, say so: ``The next two names are only landmarks; the
method appears in Section 6.2.'' A preview, a decision box, or a reference
table may legitimately name an object early, but the document should declare
that exemption rather than silently teaching the rule to ignore every early
name.

\begin{table}[H]
\centering
\small
\begin{tabularx}{\textwidth}{@{}p{0.18\textwidth}p{0.27\textwidth}p{0.25\textwidth}L@{}}
\toprule
Round & What the reader said & Located cause & Repair and product lesson \\
\midrule
1 & ``Too dense.'' & Six mechanisms and five citations in one 210-word
paragraph. & One idea per paragraph; report the passage and let the author
choose the split. \\
2 & ``Too dense.'' & New concepts arrived faster than in the named exemplars
(70 versus 83 and 253 words per concept). & Measure relative pace and burst;
use examples, figures, and read-backs before adding another term. \\
3 & ``Do not use it before explaining it.'' & About 25 first-use violations,
including an unexplained running-example algorithm. & Audit first use with a
whitelist and explicit exemption zones; introduce the generic object first. \\
\bottomrule
\end{tabularx}
\caption{One reader complaint, three levels of diagnosis. All figures are
internal calibration observations from one engagement.}
\label{tab:hpo-rounds}
\end{table}

\section{The diagnostic ladder}
\label{sec:hpo-ladder}

The case suggests four questions, in a deliberate order.

\begin{description}
\item[Surface.] Does a repeated opening, stock phrase, or private label
  deserve inspection? Pattern tools are useful here because their observations
  are narrow and easy to locate.
\item[Packing.] Does one paragraph ask the reader to hold several mechanisms,
  citations, or qualifications at once? Paragraph length is only a clue; the
  finding must show what is being carried together.
\item[Pacing.] How quickly does this document introduce and consolidate new
  concepts relative to the exemplars named in its reading brief? The meter
  reports introduction rate, singleton share, burst paragraphs, and a rolling
  live load. It does not decide whether a theorem or definition is too dense.
\item[Order.] At the first occurrence of a term, has the document supplied an
  expansion, a plain-language gloss, or a pointer to where the term is taught?
  The audit uses the reader's known-vocabulary list and declared exemption
  zones. It should abstain when the evidence is ambiguous.
\end{description}

The order matters. A writer who fixes a repeated heading before noticing that
the paragraph has no example may polish the surface while leaving the reader's
real problem untouched. Conversely, a writer who discovers an ordering fault
does not need to rewrite the whole chapter. The finding queue should be
sortable by unit, and each repaired unit should be rebuilt and inspected before
the change spreads to the rest of the document.

\section{Slow, but easy to leave}
\label{sec:hpo-skippable}

The manager offered a practical rule: when a passage is slow, the reader can
skip; when it is too fast, the reader must reconstruct the missing steps. We
take that as a repair prior, not as a licence to add padding. A slow section
earns its space when it closes questions in sequence, uses headings that make
the stopping points visible, and returns to the same idea in a different form.

The first build should record a second property beside concept pace:
skippability. A long unit with no internal heading, example, or summary is not
slow in a useful way. It is simply hard to re-enter after an interruption. The
author can respond by adding a map, a worked micro-case, or a short read-back;
the author need not accept the meter's preferred wording.

\section{The reader's evidence completes the loop}
\label{sec:hpo-reader-evidence}

The agent judged the document readable before each of the three rounds and was
wrong each time. The meters did not judge it good. They located passages that
the manager then confirmed or rejected. That division of labour belongs in the
product contract.

The reader record should preserve the moment at which the document stopped
working. In addition to a response and elapsed time, the reviewer records the
page or section where they stopped, a short quotation that triggered the
request, and the smallest missing relation they were trying to recover. Those
fields are more informative than a single score: ``I stopped at page 18,
where PPO appears without an explanation'' tells the author what to repair.

The reading brief must learn from the session. It carries the named exemplars,
the vocabulary the reader already owns, the sections in which names are only
landmarks, and the standing instruction to err on the slow side. A later
inspection runs against that versioned brief rather than reopening a settled
argument. The brief remains a reader contract, not a private project diary.

\begin{takeawaybox}{What this case adds to Humanvoice}
The HPO engagement does not prove that Humanvoice works. It gives the first
build a sharper job: distinguish packing, pacing, and ordering; calibrate
against the reader's exemplars; show the passage that triggered a finding; and
leave the repair and the final judgment with the author and reader.
\end{takeawaybox}
